\documentclass[12pt,a4paper]{article}
\usepackage{fontspec}
\usepackage[russian,english]{babel}
\babelfont{rm}{DejaVu Serif}
\usepackage[margin=2.5cm]{geometry}
\usepackage{hyperref}
\hypersetup{hidelinks}
\setlength{\parindent}{1.25cm}
\setlength{\parskip}{0.28em}
\begin{document}
\begin{center}
{\Large\bfseries Анализ бахадуровской эффективности критериев согласия}\par
\vspace{0.3cm}
{\large Материалы и обзор литературы}\par
\end{center}
\vspace{0.25cm}
\par\noindent\textbf{1. Bahadur R. R. Stochastic comparison of tests}\par
The Annals of Mathematical Statistics, 1960, vol. 31, no. 2, pp. 276–295. DOI: 10.1214/aoms/1177705894.\par
\noindent\textbf{Ссылка.} \url{https://doi.org/10.1214/aoms/1177705894}\par
\noindent\textbf{Обзор.} Работа посвящена общему подходу к стохастическому сравнению статистических критериев. Рассматриваются принципы асимптотического сопоставления процедур проверки гипотез, в том числе на основе поведения достигаемого уровня значимости.\par
Этот источник необходим для понимания исходных идей метода Бахадура и обоснования выбранного подхода к сравнению критериев согласия. На него можно опираться при изложении теоретических основ курсовой работы.\par
\par\noindent\textbf{2. Nikitin Ya. Yu. Asymptotic Efficiency of Nonparametric Tests}\par
Cambridge University Press, 1995. ISBN 0-521-47029-3.\par
\noindent\textbf{Ссылка.} \url{https://search.worldcat.org/search?q=isbn%3A0521470293}\par
\noindent\textbf{Обзор.} В монографии систематически рассматриваются подходы к асимптотической эффективности непараметрических критериев, включая точный и локальный наклоны по Бахадуру, большие уклонения и эффективность критериев согласия.\par
Монография станет основной теоретической базой для изучения условий применимости используемых формул. Она поможет корректно определить характеристики статистик, которые нужно оценивать численно.\par
\par\noindent\textbf{3. Milošević B., Nikitin Ya. Yu., Obradović M. Bahadur efficiency of EDF based normality tests when parameters are estimated}\par
arXiv:2106.07437, 2021.\par
\noindent\textbf{Ссылка.} \url{https://arxiv.org/abs/2106.07437}\par
\noindent\textbf{Обзор.} Исследуется бахадуровская эффективность критериев нормальности на основе эмпирической функции распределения при оценивании параметров. Для анализа используются приближённые наклоны, основанные на хвостовой асимптотике предельного распределения статистики и её поведении при альтернативе.\par
Это главный методологический источник для разрабатываемого алгоритма. В pysatl-experiment планируется численно оценивать коэффициент хвостовой асимптотики по результатам моделирования при нулевой гипотезе и предельное значение статистики по выборкам из альтернативных распределений.\par
\par\noindent\textbf{4. Obradović M. On asymptotic efficiency of goodness of fit tests for Pareto distribution based on characterizations}\par
Filomat, 2015, vol. 29, no. 10, pp. 2311–2324.\par
\noindent\textbf{Ссылка.} \url{https://www.pmf.ni.ac.rs/filomat/}\par
\noindent\textbf{Обзор.} В статье исследуются критерии согласия с распределением Парето, построенные на характеризациях распределения. Рассматривается их бахадуровская эффективность при различных альтернативах, что позволяет сравнивать эффективность конкретных критериев.\par
Мне нужны примеры аналитически рассчитанных эффективностей для проверки численного метода. Кроме того, работа может служить образцом для представления результатов в виде таблицы «критерий — альтернатива — эффективность».\par
\par\noindent\textbf{5. Durio A., Nikitin Ya. Yu. Local Bahadur efficiency of some goodness-of-fit tests under skew alternatives}\par
Journal of Statistical Planning and Inference, 2003, vol. 115, no. 2, pp. 171–179. DOI: 10.1016/S0378-3758(02)00155-6.\par
\noindent\textbf{Ссылка.} \url{https://doi.org/10.1016/S0378-3758(02)00155-6}\par
\noindent\textbf{Обзор.} Статья посвящена локальной бахадуровской эффективности некоторых критериев согласия при асимметричных альтернативах. Исследуется поведение критериев вблизи нулевой гипотезы.\par
Материал поможет подобрать альтернативные распределения для вычислительных экспериментов и сопоставить численные оценки с аналитическими результатами в локальном режиме. Это важно для исследования точности и сходимости метода.\par
\end{document}
